\section{Zadání}
Vybraná platonská tělesa použijte pro polyedrickou aproximaci sféry. Na plošky platonských těles znázorněte v gnomonické projekci
\begin{align*}
x &= R \cdot \tan (90^{\circ} - \textit{š}) \cdot \cos d, \\
y &= R \cdot \tan (90^{\circ} - \textit{š}) \cdot \sin d
\end{align*}
geografickou síť doplněnou zákresem kontinentů. Skript pro generování sítě poledníků, rovnoběžek a znázornění kontinentů realizujte v programu MATLAB bez použití externích funkcí. Soubory se zákresem kontinentů exportujte z vhodné datové sady.

Vytvořte prostorové modely polyedrických globů v měřítku 1:100~000~000 nebo 1:50~000~000 dle možnosti Vaší tiskárny. Součástí úlohy bude příloha s rozloženými modely globů.

Zjistěte následující vlastnosti gnomonické projekce v okrajovém bodě Q jedné ze stěn Platonského tělesa:

\begin{itemize}
    \item měřítko $m_p$ v poledníku, měřítko $m_r$ v rovnoběžce,
    \item poloosy $a$, $b$ Tissotovy indikatrix,
    \item úhel $\omega'$ mezi obrazem poledníku a rovnoběžky,
    \item maximální úhlové zkreslení $\Delta\omega$,
    \item měřítko ploch $P$.
\end{itemize}

Vypočtené parametry v bodě $Q$ použijte k zákresu Tissotovy indikatrix (doporučené měřítko 1:1 000 000), jako podklad využijte obraz geografické sítě na příslušné stěně Platonského tělesa v gnomonické projekci, volte $\Delta u = \Delta v = 10^{\circ}$ (formát A4, popř. odpovídající).

Výpočty proveďte pro referenční kouli s poloměrem $R=6380$ km, hodnoty měřítek a poloos Tissotovy indikatrix uveďte s přesností na 6 desetinných míst, úhlové hodnoty s přesností na \textit{"}. Výsledky zkontrolujte s hodnotami získanými z programu Proj.4.


